\usepackage[T1]{fontenc}
\usepackage[utf8]{inputenc}
\usepackage{amsmath,amsthm}
\usepackage{newpxtext,newpxmath}
\usepackage[margin=1.15in]{geometry}
\usepackage{mathtools}
\usepackage{microtype,booktabs,enumitem,xcolor,graphicx,bm}
\usepackage{pgfplots}
\pgfplotsset{compat=1.18}
\usepgfplotslibrary{groupplots}
\usetikzlibrary{arrows.meta,positioning,calc}
\usepackage[breakable,skins]{tcolorbox}
\usepackage[hidelinks]{hyperref}
\usepackage{fancyhdr}

\setlist{itemsep=2pt,topsep=4pt}
\setlength{\parskip}{4pt}
\setlength{\headheight}{14pt}
\pagestyle{fancy}
\fancyhf{}
\fancyhead[LE,RO]{\thepage}
\fancyhead[RE]{\small\nouppercase{\leftmark}}
\fancyhead[LO]{\small\nouppercase{\rightmark}}
\renewcommand{\headrulewidth}{0.3pt}

% ---- colours: series colours for plots, muted tints for boxes
\definecolor{ser1}{HTML}{2A78D6}
\definecolor{ser2}{HTML}{EB6834}
\definecolor{ser3}{HTML}{1BAF7A}
\definecolor{inkmuted}{HTML}{52514E}
\definecolor{gridc}{HTML}{E4E3DE}

% ---- theorem environments
\theoremstyle{plain}
\newtheorem{theorem}{Theorem}[chapter]
\newtheorem{proposition}[theorem]{Proposition}
\newtheorem{lemma}[theorem]{Lemma}
\newtheorem{corollary}[theorem]{Corollary}
\theoremstyle{definition}
\newtheorem{definition}[theorem]{Definition}
\newtheorem{example}[theorem]{Example}
\theoremstyle{remark}
\newtheorem{remark}[theorem]{Remark}

% ---- boxes
\newtcolorbox{intuition}{breakable,enhanced,colback=ser1!5,colframe=ser1!55,boxrule=0pt,leftrule=2.5pt,
  arc=0pt,outer arc=0pt,left=8pt,right=6pt,top=5pt,bottom=5pt,before skip=9pt,after skip=9pt,
  fontupper=\normalsize,title={\textbf{Intuition.}},attach title to upper={\ }, coltitle=black}
\newtcolorbox{pitfall}{breakable,enhanced,colback=ser2!6,colframe=ser2!70,boxrule=0pt,leftrule=2.5pt,
  arc=0pt,outer arc=0pt,left=8pt,right=6pt,top=5pt,bottom=5pt,before skip=9pt,after skip=9pt,
  title={\textbf{Pitfall.}},attach title to upper={\ }, coltitle=black}
\newtcolorbox{numcheck}{breakable,enhanced,colback=ser3!6,colframe=ser3!70,boxrule=0pt,leftrule=2.5pt,
  arc=0pt,outer arc=0pt,left=8pt,right=6pt,top=5pt,bottom=5pt,before skip=9pt,after skip=9pt,
  title={\textbf{Numerical check.}},attach title to upper={\ }, coltitle=black}

\newenvironment{exercises}{\section*{Exercises}\begin{enumerate}[label=\textbf{\thechapter.\arabic*},leftmargin=*]}{\end{enumerate}}

% ---- notation
\newcommand{\R}{\mathbb{R}}
\newcommand{\N}{\mathbb{N}}
\newcommand{\Z}{\mathbb{Z}}
\newcommand{\Prob}{\mathbb{P}}
\newcommand{\Q}{\mathbb{Q}}
\newcommand{\E}{\mathbb{E}}
\newcommand{\EQ}{\mathbb{E}^{\Q}}
\newcommand{\F}{\mathcal{F}}
\newcommand{\G}{\mathcal{G}}
\newcommand{\B}{\mathcal{B}}
\newcommand{\Var}{\operatorname{Var}}
\newcommand{\Cov}{\operatorname{Cov}}
\newcommand{\Corr}{\operatorname{Corr}}
\newcommand{\dd}{\mathrm{d}}
\newcommand{\ind}{\mathbf{1}}
\newcommand{\as}{\text{a.s.}}
\newcommand{\qv}[1]{\langle #1\rangle}
\newcommand{\Nc}{\mathcal{N}}
\newcommand{\sgn}{\operatorname{sgn}}
\newcommand{\VaR}{\operatorname{VaR}}
\newcommand{\ES}{\operatorname{ES}}
\newcommand{\Ind}[1]{\ind_{\{#1\}}}
\DeclareMathOperator{\tr}{tr}
\DeclareMathOperator*{\argmax}{arg\,max}
\DeclareMathOperator*{\argmin}{arg\,min}

% ---- plot style: thin lines, recessive grid and axes
\pgfplotsset{booklet/.style={
  width=0.92\linewidth,height=5.6cm,
  grid=major,major grid style={gridc,line width=0.3pt},
  axis line style={inkmuted,line width=0.4pt},tick style={inkmuted,line width=0.3pt},
  tick label style={font=\footnotesize,color=inkmuted},
  label style={font=\small},title style={font=\small},
  legend style={font=\footnotesize,draw=none,fill=none},legend cell align=left,
  every axis plot/.append style={line width=0.9pt,no markers},
}}
\newcommand{\datapath}{data}
